\section{Details of the $\knrw$ algorithm}
We here recall the main results of \cite{KNRW18}.  Let $S$ denote a
set of $n$ labeled examples $\{z_i = (x_i, x'_i), y_i)\}_{i=1}^n$, and
let $\cP$ denote the empirical distribution over this set of examples.
Let $D$ denote a probability distribution over $\cH$. Consider the
following {\em Fair ERM (Empirical Risk Minimization)} problem:
\begin{align}
  &  \min_{D\in \Delta_\cH}\; \Ex{h\sim D}{err(h, \cP)}\\
  \mbox{such that } \forall g\in \cF \qquad &
\fpsize(g, \cP) \cdot \fpdisp(g, D, \cP) \leq \gamma.
\end{align}
where $err(h, \cP) = \Pr_\cP[h(x, x') \neq y]$, and the quantities
$\fpsize$ and $\fpdisp$ are defined in \Cref{fp-fair}.
We will write $\OPT$ to denote the objective value at the optimum for
the Fair ERM problem, that is the minimum error achieved by a
$\gamma$-fair distribution over the class $\cH$. This is the fair learning problem that we want to solve.   We assume our algorithm has access to the cost-sensitive
  classication oracles $\CSC(\cH)$ and $\CSC(\cF)$ over the classes
  $\cH$ and $\cF$ respectively.

Kearns et al. \cite{KNRW18} prove the existence of an oracle-efficient algorithm for solving the Fair ERM problem:

\begin{theorem}[\cite{KNRW18}]\label{thm:polytime}
  Fix any $\nu, \delta\in (0, 1)$. Then given an input of $n$ data
  points and accuracy parameters $\nu, \delta$ and access to oracles
  $\CSC(\cH)$ and $\CSC(\cF)$, there exists an algorithm that runs in
  polynomial time, and with probability at least $1 - \delta$, outputs
  a randomized classifier $\hat D$ such that
  $err(\hat D, \cP) \leq \OPT + \nu$, and for any $g\in \cF$, the
  fairness constraint violations satisfies
  \[
    \alpha_{FP}(g, \cP) \cdot\beta_{FP}(g, \hat D, \cP) \leq \gamma + O(\nu).
  \]
\end{theorem}


The algorithm corresponding to Theorem \ref{thm:polytime} is randomized, however, and hence less amenable to the sort of empirical investigation that we undertake in this paper. Fortunately, \cite{KNRW18} also give another algorithm, with somewhat weaker guarantees. It has the same guarantees as Theorem \label{thm:polytime}, except the convergence guarantees hold only after an exponential, rather than a polynomial number of steps. However, it has the virtue of very simple per-step dynamics, and is the algorithm that we investigate in this paper. Its pseudo-code follows:

\begin{algorithm}[h]
  \caption{\bf{FairFictPlay: Fair Fictitious Play}}
 \label{alg:fairfict}
  \begin{algorithmic}
    \STATE{\textbf{Input:} distribution $\cP$ over the labelled data
      points, CSC oracles $\CSC(\cH)$ and $\CSC(\cG)$ for the classes
      $\cH(S)$ and $\cG(S)$ respectively, dual bound $C$, and number
      of rounds $T$}

    \STATE{\textbf{Initialize}: set $h^0$ to be some classifier in
      $\cH$, set $\lambda^0$ to be the zero vector. Let
      $\overline D$ and $\overline \lambda$ be the point distributions that put
      all their mass on $h^0$ and $\lambda^0$ respectively.}

    \STATE{\textbf{For} $t = 1, \ldots, T$:}

    \INDSTATE{\textbf{Compute the empirical play distributions}:}

    \INDSTATE[2]{\textbf{Let} $\overline D$ be the uniform distribution
      over the set of classifiers $\{h^0, \ldots, h^{t-1}\}$ }

    \INDSTATE[2]{\textbf{Let}
      $\overline \lambda = \frac{\sum_{t'< t} \lambda^{t'}}{t}$ be the
      auditor's empirical dual vector}


    \INDSTATE{\textbf{Learner best responds}:
      Use the oracle $\CSC(\cH)$ to compute
      $h^{t} = \argmin_{h \in \cH(S)} \langle \LC(\overline\lambda), h
      \rangle$
\iffalse
Call the oracle
      $\cC$ to find a classifier
      \[
        h^t\in \argmin_{h\in \cH} \; err(h, \cP) + \sum_{g\in
          \cG_\alpha} \overline \lambda_g \left(\FP(h, g) - \FP(h)
        \right)
      \]\fi
}

\INDSTATE{\textbf{Auditor best responds}: Use the oracle $\CSC(\cF)$
  to compute
  $\lambda^t = \argmax_{\lambda} \Ex{h\sim \overline D}{U(h,
    \lambda)}$
\iffalse
Call the oracle
    $\cA$ to find a subgroup
    \[
      g(\overline D) \in \argmax_{g\in \cG_\alpha} |\Ex{h\sim \overline
        p}{\FP(h, g)} - \Ex{h\sim \overline D}{\FP(h)}|
    \]
\fi
}
\iffalse  \INDSTATE{
    Let $\lambda^t \in \Lambda_{\text{pure}}$ with
    $\lambda^t_{g(\overline D)} = C \; \text{sgn}\left[ \Ex{h\sim
        \overline D}{\FP(h, g(\overline D))} - \Ex{h\sim \overline
        p}{\FP(h)}\right]$}\fi


    \STATE{\textbf{Output:} the final empirical distribution
      $\overline D$ over classifiers}

    \end{algorithmic}
  \end{algorithm} 
  
  To briefly introduce the notations in the description above, we
  first note that we can rewrite the set of constraints in the Fair
  ERM problem as follows: for each $g\in \cG(S)$, 
\begin{align}
    &\Phi_+(h, g) \equiv \fpsize(g, P) \, \left(\FP(h) - \FP(h, g) \right) - \gamma\leq 0 \label{f1}\\
  &  \Phi_-(h, g)\equiv \fpsize(g, \cP)\, \left(\FP(h, g) - \FP(h)\right) - \gamma \leq 0 \label{f2}
\end{align}
Here $\cG(S)$ and $\cH(S)$ denote the set of all labellings on $S$
that are induced by $\cG$ and $\cH$ respectively, that is
\begin{align}
  &\cG(S) = \{(g(x_1), \ldots , g(x_n)) \mid g\in \cG\} \qquad
  \mbox{and,} \\
&\cH(S) = \{(h(X_1), \ldots , h(X_n))\mid h\in
  \cH\}
\end{align}
Then, $\lambda$ is a vector with a coordinate $\lambda_g^+$ and
$\lambda_g^-$ for every subgroup $g \in \cF$ such that $\lambda_g^+$
and $\lambda_g^-$ are dual variables that corresponds the pair of
constraints \eqref{f1} and \eqref{f2}.  The partial Lagrangian of the
linear program is the following:
\begin{align*}
  \cL(D, \lambda) = \Ex{h\sim D}{err(h, \cP)} +  \sum_{g\in \cG(S)}
  \left( \lambda_g^+\,  \Phi_+(D, g) +  \lambda_g^- \,  \Phi_-(D,g) \right)
\end{align*}
Similarly, the payoff function for the zero-sum game is then defined
as: for any pair of actions
$(h, \lambda)\in \cH\times \Lambda_{\text{pure}}$,
\[
  U(h, \lambda) = err(h, \cP) + \sum_{g\in \cG(S)}
  \left(\lambda_g^+\Phi_+(h, g) + \lambda_g^-\Phi_-(h, g) \right).
\]
Given a fixed set of dual variables $\lambda$, we will write
$\LC(\lambda) \in \RR^n$ to denote the vector of costs for labelling
each datapoint as $1$. That is, $\LC(\lambda)$ is the vector such that
for any $i \in [n]$, $\LC(\lambda)_i = c_i^1$.

We can find a best response for the Learner by making a call to the
cost-sensitive classification oracle. In particular, we assign costs
to each example $(X_i, y_i)$ as follows:
\begin{itemize}
% \item if $y_i = 1$, then $c_i^0 = \frac{1}{n}$ and $c_i^1 = 0$;
\item if $y_i = 1$, then $c_i^0 = 0$ and $c_i^1 = - \frac{1}{n}$;
\item otherwise, $c_i^0 = 0$ and
  \begin{align}
    c_i^1 = \frac{1}{n} &+ \frac{1}{n}\sum_{g\in \cG(S)}
                          (\lambda_g^+ - \lambda_g^-) \left(\Pr[g(x) = 1\mid y = 0] -  \mathbf{1}[g(x_i) = 1]\right)
  \end{align}
\end{itemize}
Then given a fixed set of dual variables $\lambda$, we will write
$\LC(\lambda) \in \RR^n$ to denote the vector of costs for labelling
each datapoint as $1$. That is, $\LC(\lambda)$ is the vector such that
for any $i \in [n]$, $\LC(\lambda)_i = c_i^1$.
